\documentclass[11pt]{article}

\usepackage[a4paper,margin=1in]{geometry}
\usepackage{amsmath,amssymb}
\usepackage[hidelinks]{hyperref}

\newcommand{\BM}{\mathcal{BM}}
\newcommand{\dbm}{d_{\mathrm{BM}}}

\title{Literature Status: Higher-Dimensional Equilateral Sets In Banach--Mazur Compacta}
\author{Automatic math-discovery smoke test}
\date{July 19, 2026}

\begin{document}
\maketitle

\section*{Status}

\textbf{Already answered in the literature}.  The higher-dimensional
equilateral-set existence question raised in the source paper is answered by a
separate later paper.

\section*{Source Question}

The source paper is Tomasz Kobos and Konrad Swanepoel,
\emph{Equilateral dimension of the planar Banach--Mazur compactum},
arXiv:2503.21478 \cite{KobosSwanepoel2025}.  The relevant passage appears in
the Concluding remarks on PDF page 13 and source lines 424--426.

In compact form, the source passage says that it is not clear whether the
planar construction can be generalized to higher dimensions.  It then states
that most likely the \(n\)-dimensional Banach--Mazur compactum contains
arbitrarily large equilateral sets for every \(n\ge2\).

\section*{Supporting Answer}

The supporting paper is Florian Grundbacher and Tomasz Kobos,
\emph{Exact Banach--Mazur distances of certain \(\ell_p\)-sums and cones},
arXiv:2603.18268 \cite{GrundbacherKobos2026}.  The authors explicitly connect
their result to the source paper: the abstract says that they lift a recent
construction of arbitrarily large equilateral sets in the \(2\)-dimensional
symmetric compactum to all higher dimensions.

The decisive statement is Corollary 1.9 on PDF page 5:
\[
  \text{for every } n,m\ge2,\text{ there are equilateral sets of cardinality }
  m \text{ in } \BM_s^n.
\]
This directly answers the higher-dimensional existence question in the
symmetric Banach--Mazur compactum.  The proof of Corollary 1.9 appears near the
end of Section 4, on PDF page 23.  It applies the local isometric embedding of
Corollary 1.8 to the equilateral family from Kobos--Swanepoel.

\section*{Scope Limitations}

This status packet records an already-known answer to the source's
higher-dimensional existence question.  It does not resolve the source paper's
separate planar question about obtaining equilateral sets for every distance
in an interval \((1,c]\), nor the broader speculation that the Banach--Mazur
compactum might contain all finite metric spaces after scaling distances close
to \(1\).

\section*{Search Evidence}

The local lightweight indexes
\[
\begin{gathered}
\texttt{registry\_index.tsv},\quad
\texttt{solutions/index.tsv},\\
\texttt{attempts/index.tsv},\quad
\texttt{proof\_gaps/index.tsv}
\end{gathered}
\]
were searched for arXiv:2503.21478, the paper title, ``equilateral
dimension'', ``planar Banach--Mazur'', ``Banach--Mazur compactum'', and
``equilateral''.  The hits were broader equilateral-set packets and attempts,
but no duplicate packet for this exact higher-dimensional source question.

A local parsed-source search found arXiv:2603.18268.  Its abstract,
Corollary 1.9, and proof of Corollary 1.9 explicitly cite/use the planar
Kobos--Swanepoel construction and lift it to all dimensions.

\section*{Packet Files}

The original source PDF is stored as \texttt{source\_paper.pdf}.  The decisive
supporting answer PDF is stored as
\texttt{supporting\_paper\_2603.18268.pdf}.

\begin{thebibliography}{9}

\bibitem{KobosSwanepoel2025}
T.~Kobos and K.~Swanepoel,
\emph{Equilateral dimension of the planar Banach--Mazur compactum},
arXiv:2503.21478; Proc. Amer. Math. Soc. 153(10) (2025), 4423--4436.

\bibitem{GrundbacherKobos2026}
F.~Grundbacher and T.~Kobos,
\emph{Exact Banach--Mazur distances of certain \(\ell_p\)-sums and cones},
arXiv:2603.18268, 2026.

\end{thebibliography}

\end{document}
